\documentclass[11pt]{article}
\input{../MilestoneDraft.text}
\input{../Common.text}
\title{Problem Statement and Goals\\\progname}
\author{\authname}
\date{September 23, 2026}
\begin{document}
\maketitle
\draftnotice

\begin{table}[h]
\centering
\caption{Revision history.}\label{tab:ps-revisions}
\begin{tabularx}{\textwidth}{llX}
\toprule
Date & Revision & Change \\
\midrule
2026-09-23 & Review draft & AI-assisted draft from Aaron Loh's decisions;
whole-team review pending. \\
\bottomrule
\end{tabularx}
\end{table}
Table~\ref{tab:ps-revisions} records this draft's origin. Natural Control is
a working name. The accompanying Development Plan contains the proof of
concept (POC) plan and proposed team charter.

\section{Problem Statement}
Blind computer users can encounter barriers when completing tasks across
websites, applications, and connected services. The project addresses
independent task completion and the ability to remain informed and in
control when delegating digital work.

\subsection{Problem}
Completing a digital task can require navigating unfamiliar interfaces,
tracking changes, and recovering from controls or content that existing
access tools do not expose effectively. Screen readers already provide
selective navigation; users do not simply listen to every item on a page.
WebAIM's self-selected survey identifies barriers including unclear
controls and unexpected changes, but does not establish their prevalence
among our prospective participants \citep{WebAIM2024}.

Delegation to an artificial intelligence (AI) agent introduces another
problem: the user needs to understand what is happening, make meaningful
choices, and correct mistakes without visual monitoring. Research on Morae
examines this issue by involving blind and low-vision users in decisions
during agent execution \citep{Morae2025}. These sources motivate the
project; interviews must still establish participants' actual priorities.

Natural Control aims to make general digital task delegation usable
through conversation while preserving user agency. Success means a person
can complete and verify useful tasks independently after initial setup,
not that they abandon their screen reader. Representative tasks span
connected accounts, unfamiliar websites, and local computer applications.
They are evaluation cases, not a product allowlist or fixed scripts.

\subsection{Inputs and Outputs}
The intended inputs are spoken goals, preferences, clarification answers,
corrections, approvals, and stop requests. Users may also supply keyboard
or text input. With appropriate access, the system receives relevant page,
application, file, and connected-account state needed for the requested
task. Initial setup supplies permitted accounts, tools, and provider
credentials; credentials are not intended for spoken entry or public logs.

The intended outputs are actions on permitted services or the computer;
concise spoken progress and results; requests for clarification or
approval; and an accessible, reviewable task history. Completion feedback
must distinguish verified results, partial work, and unresolved failures.
An acknowledgment of a stop request is not a claim that execution has
already stopped or that prior actions were undone.

\subsection{Stakeholders}
\begin{description}[style=nextline]
\item[Primary users.]
Blind people with screen-reader experience who want to delegate digital
tasks while retaining control. Initial evaluation is with English-speaking
Mac users; this is a development boundary, not a claim that these users
represent all blind people.
\item[Research collaborators and accessibility organizations.]
Potential participants and advisers can identify barriers and evaluate
interaction choices. A team connection may facilitate a CNIB introduction;
no partnership, participation, or endorsement is confirmed.
\item[Setup helpers and open-source contributors.]
Trusted helpers may assist with initial installation and account setup.
Maintainers need understandable code, documented limitations, and a
reproducible default configuration.
\item[Course team and assessors.]
Four student contributors, the teaching assistant (TA), and instructor
need a feasible project and evidence that the team understands its work.
People receiving messages or invitations are also affected by the
system's actions, which makes accurate user authorization important.
\end{description}

\subsection{Environment and Boundaries}
The initial evaluated environment is macOS with Chrome, English speech,
a microphone, audio output, internet access, and existing accessibility
tools such as VoiceOver and keyboard navigation. A web voice interface is
planned; its location in a browser does not restrict the agent to web
tasks. Cloud processing is permitted, and users provide their own account
connections and provider credentials.

A local component is necessary for local Mac tasks. An always-available
Mac or Mac mini is one deployment option; hosted execution is another for
tasks that do not require a local device. A hosted service cannot operate
an offline Mac. Independent everyday use is the initial goal; assisted
setup is acceptable. Telephone access will be compared with on-Mac voice,
with support for both channels treated as an extension.

Unrestricted success on every website or application, guaranteed
reversibility, universal provider compatibility, and replacement of all
screen-reader functions are not promised. The project will build on an
existing agent rather than train a new foundation model. User needs and
provider constraints may change the evaluated scope, with changes
recorded rather than hidden.

\section{Goals}
The following four goals describe benefits to users. Numerical targets
are provisional design commitments accepted for this draft, not measured
results. A pilot will inform revisions before the main evaluation; changes
will be justified and retained in the revision history. Detailed task
selection and measurement procedures belong in the Development Plan and
later verification and validation plan.

\begin{enumerate}[label=G\arabic*.,leftmargin=*]
\item \textbf{Complete useful digital tasks independently.}
After assisted setup, users should complete at least 80\% of a declared,
representative task set correctly without sighted assistance. The set will
cover connected-account, browser, and local-Mac tasks. Report results per
category and overall, including failed attempts. This target applies to
the evaluated tasks, not arbitrary computer use.

\item \textbf{Retain control over delegated work.}
Users should be able to stop work and decide whether consequential
actions proceed. In the evaluation, explicit stop requests should be
acknowledged within two seconds, and every tested consequential action
must require explicit approval. Verified stopping is separately checked:
no new action may be issued after the system confirms execution stopped.

\item \textbf{Understand and correct the assistant without visual help.}
In each evaluated scenario, check whether the participant can identify the
actual outcome or blocker and, in a seeded-misunderstanding scenario,
correct the request without sighted assistance. Report successful cases
and failure reasons. Brief feedback and access to details should support
awareness without requiring constant narration.

\item \textbf{Make the system usable beyond the development team.}
Release an open-source implementation with one documented default
configuration, user-provided credentials, and extension points. Verify
that a person outside the implementation team can follow the setup guide
on a clean supported environment, with assistance documented, and run the
reference tasks. Stars or downloads alone are not evidence of usefulness.
\end{enumerate}

\section{Stretch Goals}
The leading extension is a second voice channel: let a user access the
same task context through telephone and on-Mac voice. Later possibilities
include messaging continuity and evaluated Safari support. These must not
displace a reliable initial experience or be presented as completed work.

\section{Custom Documents (CDs)}
Two custom documents are proposed under the course outline
\citep{CourseOutline}. \textbf{Instructor approval is pending.}
\begin{itemize}
\item \textbf{Fall: User/Stakeholder Interview Report.} Record recruitment,
participant needs, methods, findings, and resulting design decisions.
Separate direct evidence from team hypotheses.
\item \textbf{Winter: Usability Report.} Evaluate independent completion,
awareness, correction, and user control against participants' usual
workflows. Report limitations and unsuccessful cases as well as benefits.
\end{itemize}

\section*{Appendix A: Generative AI Use Disclosure}
\addcontentsline{toc}{section}{Appendix A: Generative AI Use Disclosure}
\textbf{Tool and use.} Aaron Loh used OpenAI Codex with GPT-6 Astra,
primarily at Medium reasoning and occasionally High (settings reported by
Aaron). Codex helped extract questions from course requirements, challenge
assumptions, research options, and develop ideas through discussion. It also
assisted with drafting and revising text throughout this document, organizing
references, and LaTeX/PDF formatting. Aaron made the judgment calls and reviewed
the work. These are working drafts for further team revision, not evidence of
team-wide agreement or completed participant research.

\textbf{Verification and responsibility.} The assistant checked cited sources,
course requirements, document consistency, and PDF rendering. These checks do
not establish that the proposed system works. Team review, ratification, and
personal reflections remain pending. Before submission, the team will verify
the final content and update this disclosure to reflect subsequent revisions
and any additional AI use. The team remains responsible for the submitted work.

\section*{Appendix B: Reflection}
These prompts require the team's own account of producing the deliverable.
They have deliberately not been answered on behalf of absent contributors.
\begin{enumerate}
\item \textbf{What went well?} \pending{Describe actual collaboration and
what helped produce this deliverable.}
\item \textbf{What pain points arose and how were they resolved?}
\pending{Supply actual experiences and actions, not hypothetical lessons.}
\item \textbf{How did the team adjust scope?} \pending{Reflect on the
scope decision. The planning record establishes a general-agent interface
with bounded evaluation, but does not establish every member's reasoning
or agreement.}
\end{enumerate}

\clearpage
\bibliographystyle{plainnat}
\bibliography{../../refs/CapstoneSources}
\end{document}
